% © 2026 Ing. David Jaroš — CC BY-NC-ND 4.0
%
% This work is licensed under a Creative Commons Attribution-NonCommercial-NoDerivatives
% 4.0 International License (CC BY-NC-ND 4.0).
%
% License History: Earlier drafts (up to v0.3) were released under CC BY 4.0.
% From v0.4 onward, all material is released under CC BY-NC-ND 4.0 to protect
% the integrity of the theoretical work during ongoing academic development.

\documentclass[11pt,a4paper]{article}
\usepackage{amsmath, amssymb, amsthm}
\usepackage{hyperref}
\usepackage{geometry}
\geometry{margin=2.5cm}

\newtheorem{theorem}{Theorem}
\newtheorem{proposition}{Proposition}
\newtheorem{definition}{Definition}
\newtheorem{remark}{Remark}
\newtheorem{conjecture}{Conjecture}

\title{Fermion Emergence from the \texorpdfstring{$\Theta$}{Theta}-Field\\
\large Spinor Representations, Dirac Limit, and Yukawa Couplings in UBT}
\author{Ing.\ David Jaroš}
\date{2026}

\begin{document}
\maketitle

\begin{abstract}
We investigate whether fermionic degrees of freedom emerge naturally from the
biquaternionic $\Theta$-field of Unified Biquaternion Theory (UBT).  Three
complementary approaches are examined: (1) the algebraic identification of
spinor representations within the biquaternion algebra
$\mathbb{B} = \mathbb{H}\otimes\mathbb{C} \cong \mathrm{Mat}(2,\mathbb{C})$;
(2) the derivation of the Dirac equation as an effective limit of the UBT
master equation; and (3) the topological mechanism through which
odd-winding KK modes acquire fermionic statistics.
Yukawa coupling candidates arising from the cubic term in the UBT action
are catalogued.  All claims are labelled with their proof status.
\end{abstract}

\tableofcontents

% ======================================================================
\section{The Biquaternion Algebra as a Spinor Algebra}
\label{sec:algebra}
% ======================================================================

\subsection{Algebraic isomorphism}

The biquaternion algebra $\mathbb{B} = \mathbb{H}\otimes_\mathbb{R}\mathbb{C}$
satisfies (see \texttt{canonical/fields/theta\_field.tex}, §Matrix Structure):
\begin{equation}
\label{eq:iso}
  \mathbb{B} \;\cong\; \mathrm{Mat}(2,\mathbb{C}).
\end{equation}

This is the \emph{same} algebra as the one generated by the $2\times 2$
Pauli matrices $\{\sigma_0, \sigma_1, \sigma_2, \sigma_3\}$.

\begin{proposition}[Gamma matrices from biquaternions]\label{prop:gamma}
  The Dirac gamma matrices $\gamma^\mu$ satisfying
  $\{\gamma^\mu,\gamma^\nu\} = 2\eta^{\mu\nu}$
  can be realised within $\mathbb{B}\otimes\mathbb{B}$ as:
  \begin{equation}
    \gamma^0 = \sigma_3\otimes\sigma_0, \quad
    \gamma^i = i\sigma_1\otimes\sigma_i \quad (i=1,2,3).
  \end{equation}
  \textbf{Status: [L0]} — algebraic identity.  Source:
  \texttt{consolidation\_project/qed\_phi\_const/step2\_electron\_mass.tex} §2.1.
\end{proposition}

\subsection{Spinor subspace of \texorpdfstring{$\mathbb{B}$}{B}}

\begin{definition}[Spinor subspace]
  The spinor subspace $\mathcal{S} \subset \mathbb{B}\otimes\mathbb{B}$ is
  the minimal left ideal:
  \begin{equation}
    \mathcal{S} = (\mathbb{B}\otimes\mathbb{B})\; P_+,
    \qquad P_+ = \tfrac{1}{2}(1 + \gamma^0).
  \end{equation}
  A Dirac spinor $\psi \in \mathcal{S}$ arises as the restriction of
  $\Theta$ to this left ideal.
\end{definition}

\begin{remark}[Status]
  This algebraic construction follows the standard ideal approach to
  spinors in Clifford algebras, applied here to $\mathbb{B}\otimes\mathbb{B}$.
  Status: \textbf{[L1]} — proved given the isomorphism \eqref{eq:iso}
  and the choice of left ideal.
\end{remark}

% ======================================================================
\section{Dirac Equation as an Effective Limit}
\label{sec:dirac}
% ======================================================================

\subsection{UBT master equation}

The UBT field equation (\texttt{canonical/fields/theta\_field.tex}) is:
\begin{equation}
\label{eq:master}
  \nabla^\dagger\nabla\,\Theta(q,\tau) = \kappa\,\mathcal{T}(q,\tau),
\end{equation}
where $\nabla^\dagger$ is the biquaternionic conjugate derivative,
$\mathcal{T}$ is the energy-momentum-current source tensor, and
$\kappa = 8\pi G$ in the gravitational limit.

\subsection{Real-sector projection and Dirac limit}

Setting $\psi = 0$ (real-time limit) and projecting onto the spinor subspace
$\mathcal{S}$ defined in the previous section:

\begin{theorem}[Dirac equation from UBT]\label{thm:dirac}
  In the limit $\psi \to 0$, $A_\mu \neq 0$ (electroweak sector), and
  $\Theta \in \mathcal{S}$, the master equation \eqref{eq:master} reduces to:
  \begin{equation}
  \label{eq:dirac}
    (i\gamma^\mu D_\mu - m)\,\psi = 0,
  \end{equation}
  where $D_\mu = \partial_\mu + igA_\mu$ is the U(1) gauge-covariant derivative
  and $m$ is the effective mass parameter.
  \textbf{Status: [L1]}.
  Source: \texttt{canonical/bridges/QED\_complete\_chain.tex} Theorem D.4 and
  \texttt{consolidation\_project/FPE\_verification/step3\_dirac\_emergence.tex}.
\end{theorem}

\begin{remark}[Phase-curvature corrections]
  The full biquaternionic equation \eqref{eq:master} contains additional
  phase-curvature terms $\sim \partial_\psi\Theta$ which vanish in the
  $\psi \to 0$ limit.  At finite $R_\psi$, they represent higher-dimensional
  corrections to the Dirac equation suppressed by $m/R_\psi^{-1}$.
\end{remark}

% ======================================================================
\section{Fermionic Statistics from Winding Parity}
\label{sec:winding}
% ======================================================================

\subsection{The \texorpdfstring{$\mathbb{Z}_2$}{Z2} grading from \texorpdfstring{$\psi$}{psi}-winding}

The KK mode expansion (Section 4 of \texttt{research/theta\_quantum\_field.tex}) yields
modes $\hat\Theta_n$ labelled by $n \in \mathbb{Z}$.

\begin{definition}[$\psi$-parity]\label{def:psi_parity}
  Define the $\psi$-parity operator:
  \begin{equation}
    \mathcal{P}_\psi\colon \hat\Theta_n \mapsto (-1)^n \hat\Theta_n.
  \end{equation}
  Even-winding modes ($n$ even) have $\mathcal{P}_\psi = +1$; they are
  \emph{bosonic}.  Odd-winding modes ($n$ odd) have $\mathcal{P}_\psi = -1$;
  they are candidates for \emph{fermionic} degrees of freedom.
\end{definition}

\begin{remark}[Status]
  The $\psi$-parity grading is a well-defined algebraic operation following
  from the KK expansion.  Status: \textbf{[L0]}.
  
  Whether odd-winding modes obey \emph{anticommutation relations} (and hence
  full fermionic statistics) requires a more detailed analysis.  The
  necessary condition is that the path-integral measure be Grassmannian for
  these modes.  Status: \textbf{[L2]} — motivated conjecture, not yet proved.
  Relating this to spin-statistics requires identifying the spin of
  odd-winding modes from the Θ algebra.
\end{remark}

\subsection{Three generations from winding modes}

\begin{proposition}[Three-generation structure]\label{prop:gen}
  The three lowest winding modes $n = 1, 2, 3$ provide three independent
  copies of the spinor field $\psi$, corresponding to the three fermion
  generations.
  \textbf{Status: [L0]} — proved algebraic identity.
  Source: \texttt{DERIVATION\_INDEX.md}: ``$\psi$-modes as independent
  B-fields [L0]: Proven''.
\end{proposition}

\begin{remark}[Generation mass hierarchy]
  While the three-generation structure is proved (Proposition~\ref{prop:gen}),
  the mass \emph{ratios} between generations are not derived from first
  principles.  See \texttt{research/generation\_structure.tex} and
  \texttt{research/fermion\_mass\_generation.tex}.
\end{remark}

% ======================================================================
\section{Topological Excitations and Hopfion Sector}
\label{sec:topological}
% ======================================================================

\subsection{Hopfion solutions}

Beyond the perturbative KK spectrum, the UBT field equation admits
topological soliton solutions.  In the imaginary sector, these appear as
Hopf fibrations:
\begin{equation}
  \pi_3(S^2) = \mathbb{Z},
\end{equation}
supporting stable particle-like solutions with integer topological charge
(Hopf invariant $H \in \mathbb{Z}$).

\begin{remark}[Status]
  The existence of Hopf soliton solutions in the UBT scalar sector is
  motivated by the fibration structure of $\mathbb{H}\otimes\mathbb{C}$.
  Status: \textbf{[L2]} — motivated conjecture with supporting calculation in
  \texttt{unified\_biquaternion\_theory/solution\_P5\_dark\_matter/ThetaM\_MassHierarchy.tex}.
\end{remark}

\subsection{Spin from topology}

For a Hopfion of Hopf charge $H$, the spin angular momentum takes the
value:
\begin{equation}
  J = \tfrac{1}{2} H \pmod{1}.
\end{equation}
Odd-charge Hopfions ($H$ odd) would therefore carry half-integer spin,
consistent with fermionic statistics.

\begin{remark}[Status: Conjectural]
  The identification of Hopf charge with spin-$1/2$ is a motivated
  conjecture analogous to the Skyrme-baryon connection in nuclear physics.
  A rigorous derivation requires analysing the angular-momentum eigenvalues
  of the biquaternionic soliton solutions.
  Status: \textbf{[L2]}.
\end{remark}

% ======================================================================
\section{Yukawa Couplings}
\label{sec:yukawa}
% ======================================================================

\subsection{Origin from the cubic action term}

The UBT action contains a cubic interaction between the scalar and spinor
components of $\Theta$:
\begin{equation}
\label{eq:yukawa_action}
  S_{\mathrm{Yukawa}}
  = -y\int d^4x\, d\tau\;
    \overline{\psi}(x,\tau)\, \phi(x,\tau)\, \chi(x,\tau),
\end{equation}
where $\phi$ is the scalar (Higgs) component, $\psi, \chi$ are left- and
right-handed spinor components, and $y$ is the Yukawa coupling.

\begin{remark}[Origin]
  This term arises from the cubic coupling in $V(\Theta)$ when decomposed
  in the spinor-scalar basis.  The form is identical to Standard Model
  Yukawa couplings.  Source: \texttt{consolidation\_project/masses/yukawa\_structure.tex}.
  Status: \textbf{[L1]}.
\end{remark}

\subsection{Yukawa matrix and Higgs holonomy}

Following \texttt{unified\_biquaternion\_theory/fermion\_mass\_derivation\_complete.tex}
§Yukawa, the Yukawa coupling matrix in generation space takes the form:
\begin{equation}
  Y_{ij} = y_0\, \Phi_{ij},
\end{equation}
where $y_0$ is an overall scale and $\Phi_{ij}$ is the holonomy profile on
the internal torus $\mathbb{T}^2$ (spanned by the imaginary directions
$\chi, \xi$).

\begin{remark}[Higgs as \texorpdfstring{$\Theta$}{Theta} component]
  The Higgs field $\phi$ is the zero-mode ($n=0$) of the scalar component
  of $\Theta$ in the KK expansion.  The Yukawa coupling $y_0$ is therefore
  the effective biquaternionic overlap integral:
  $y_0 = \int_{T^2} \hat\Theta_0(\chi,\xi)\, d\chi\, d\xi$.
  Its value depends on $R_\psi$ and is not yet derived from first principles.
  Status: \textbf{[O]}.
\end{remark}

% ======================================================================
\section{Research Tracks and Open Problems}
\label{sec:open}
% ======================================================================

\begin{enumerate}
  \item \textbf{[O] Grassmannian path-integral measure for odd-$n$ modes}:
    Proving that the Θ path integral for odd-winding modes naturally leads
    to anticommuting variables (Grassmann algebra) rather than commuting ones
    requires analysing the Jacobian of the mode change of variables.

  \item \textbf{[O] Hopfion spin quantisation}:
    A full derivation of fermionic spin for odd-charge Hopfions from the
    biquaternionic Noether theorem.

  \item \textbf{[O] Yukawa coupling from moduli}:
    Deriving $Y_{ij}$ entirely from $R_\psi$ and the modular parameter
    $\tau_{\mathrm{mod}}$ without fitting.

  \item \textbf{[O] Chirality selection}:
    Explaining why only $\mathrm{SU}(2)_L$ (not $\mathrm{SU}(2)_R$) couples
    to fermions.  Status: motivated by $\psi$-parity but not proved.
    Source: \texttt{canonical/bridges/gauge\_emergence\_bridge.tex} §Chirality.
\end{enumerate}

% ======================================================================
\section{Cross-References}
\label{sec:xref}
% ======================================================================

\begin{tabular}{ll}
\hline
Topic & File \\
\hline
Gamma matrices from biquaternions & \texttt{consolidation\_project/qed\_phi\_const/step2\_electron\_mass.tex} \\
Dirac emergence (Step 3) & \texttt{consolidation\_project/FPE\_verification/step3\_dirac\_emergence.tex} \\
QED chain (Dirac eq.) & \texttt{canonical/bridges/QED\_complete\_chain.tex} \\
Three-generation proof & \texttt{DERIVATION\_INDEX.md} \\
Yukawa structure & \texttt{consolidation\_project/masses/yukawa\_structure.tex} \\
Fermion mass program & \texttt{research/fermion\_mass\_program.md} \\
Hopfion masses & \texttt{unified\_biquaternion\_theory/solution\_P5\_dark\_matter/ThetaM\_MassHierarchy.tex} \\
Chirality & \texttt{canonical/bridges/gauge\_emergence\_bridge.tex} \\
\hline
\end{tabular}

\end{document}
